
\subsection{交流与双棒}

\begin{tabularx}{\columnwidth}{XX}
  磁通量&$\Phi=BS\cos\theta$\\
  线圈转动电动势&$e=nBS\omega\sin\theta$\\
  最大电动势&$E_m=nBS\omega$\\
  从中性面计时&$e=E_m\sin\omega t$\\
  周期与角速度&$T=\dfrac{2\pi}{\omega}$\\
  双棒动量关系&$m_1v_1+m_2v_2=\text{常量}$（无外力且安培力等大反向）\\
\end{tabularx}

\begin{formulanotebox}
中性面是线圈平面与磁感线垂直的位置。过中性面时磁通量最大，但磁通量变化率为零，所以感应电动势为零。

线圈平面与磁感线平行时磁通量为零，但磁通量变化最快，感应电动势最大。
\end{formulanotebox}

\begin{keypointbox}[title={\strut\textcolor{white}{\bfseries 重点知识点：交流电产生}}]
线圈匀速转动时，感应电动势按正弦规律变化。每经过一次中性面，线圈切割磁感线方向整体反向，感应电动势方向改变一次。

若从中性面开始计时，电动势先从零增大到最大值，再减小到零，随后反向重复。
\end{keypointbox}

\begin{keypointbox}[title={\strut\textcolor{white}{\bfseries 重点知识点：双棒模型}}]
\textbf{安培力等大反向}：两棒处于同一回路且 $BL$ 相同，安培力常可看成内力。无外力时整体动量守恒，最终常趋于共速，减少的机械能转化为焦耳热。\\

\textbf{安培力不等}：若两棒所在磁场、有效长度或电流关系不同，整体分析时不能直接消去安培力，应保留安培力总冲量和总功。\\

\textbf{外力作用}：有外力牵引一根棒时，对整体看外力冲量和外力做功，对单棒再分别分析安培力。
\end{keypointbox}

\begin{casetypebox}[title={\strut\textcolor{white}{\bfseries 常见题型：线圈转动与双棒综合}}]
\begin{itemize}[leftmargin=1.5em]
  \item 线圈转动题先找中性面，再判断电动势零值、最大值和方向变化。
  \item 求瞬时值时明确计时起点；从中性面开始通常用 $e=E_m\sin\omega t$。
  \item 双棒无外力且最终共速时，先用动量守恒求共同速度，再用能量守恒求焦耳热。
  \item 含外力双棒题常用整体动量定理求速度关系，用整体功能关系求热量。
\end{itemize}
\end{casetypebox}
